%%%%%%%%%%%%%%%%%%%%%%%%%%%%%%%%%%%%
%% PREAMBLE
%%%%%%%%%%%%%%%%%%%%%%%%%%%%%%%%%%%%

\documentclass[11pt,a4paper]{article}

% Taal en typografie
\usepackage[utf8]{inputenc}
\usepackage[T1]{fontenc}
\usepackage[dutch]{babel}
\usepackage{lmodern}
\usepackage{microtype}

% Wiskunde en tabellen
\usepackage{amsmath,amssymb}
\usepackage{booktabs}
\usepackage{array}
\usepackage{longtable}

% Bomen en figuren
\usepackage{graphicx}
\usepackage{qtree}

% Pagina-opmaak
\usepackage[margin=2.5cm]{geometry}
\setlength{\parindent}{0pt}
\setlength{\parskip}{6pt}

% Links
\usepackage[hidelinks]{hyperref}

% Kleine hulpfuncties voor consistente opmaak
\newcommand{\nt}[1]{\ensuremath{\mathrm{#1}}}
\newcommand{\lex}[1]{\textit{#1}}
\newcommand{\ruleprob}[2]{\ensuremath{P(#1\rightarrow #2)}}

\begin{document}

%%%%%%%%%%%%%%%%%%%%%%%%%%%%%%%%%%%%
%% HEADING
%%%%%%%%%%%%%%%%%%%%%%%%%%%%%%%%%%%%

\rule{\textwidth}{1pt}
\begin{center}
  \textbf{Taaltheorie en Taalverwerking -- Homework Submission}
\end{center}
\rule{\textwidth}{0.5pt}

%%%%%%%%%%%%%%%%%%%%%%%%%%%%%%%%%%%%
%% ENTER DETAILS
%%%%%%%%%%%%%%%%%%%%%%%%%%%%%%%%%%%%

\textbf{Week:} 3

\textbf{Group:} 28

\textbf{Student names and IDs:} Bob Engel 12627143, Dion Streefkerk 15847195

\textbf{Date:} 4 mei, 2026

\rule{\textwidth}{1pt}

%%%%%%%%%%%%%%%%%%%%%%%%%%%%%%%%%%%%
%% YOUR ANSWERS
%%%%%%%%%%%%%%%%%%%%%%%%%%%%%%%%%%%%

\begin{enumerate}

%------------------------------------------------------------
\item \textbf{CFG voor de drie CoreNLP-bomen}

De vier parameters van de grammatica zijn $\langle \Sigma, N, \nt{S}, R\rangle$.

\[
\begin{aligned}
\Sigma = \{&\text{the}, \text{a}, \text{architect}, \text{obtained}, \text{sketchbook}, \text{or}, \text{pencil}, \text{from}, \text{landlord},\\
&\text{drew}, \text{sketch}, \text{of}, \text{house}, \text{and}, \text{garden}, \text{blueprint}, \text{museum}, \text{donation}\}.
\end{aligned}
\]

\[
N = \{\nt{S}, \nt{NP}, \nt{VP}, \nt{PP}, \nt{DT}, \nt{NN}, \nt{VBD}, \nt{IN}, \nt{CC}\}.
\]

Het startsymbool is $\nt{S}$.

De regels $R$ zijn:

\[
\begin{array}{rcl}
\nt{S} &\rightarrow& \nt{NP}\ \nt{VP} \\
\nt{VP} &\rightarrow& \nt{VBD}\ \nt{NP} \mid \nt{VBD}\ \nt{NP}\ \nt{PP} \\
\nt{NP} &\rightarrow& \nt{DT}\ \nt{NN} \mid \nt{NP}\ \nt{PP} \mid \nt{NP}\ \nt{CC}\ \nt{NP} \\
\nt{PP} &\rightarrow& \nt{IN}\ \nt{NP} \\[1mm]
\nt{DT} &\rightarrow& \text{the} \mid \text{a} \\
\nt{NN} &\rightarrow& \text{architect} \mid \text{sketchbook} \mid \text{pencil} \mid \text{landlord} \mid \text{sketch} \mid \text{house} \\
   &\mid& \text{garden} \mid \text{blueprint} \mid \text{museum} \mid \text{donation} \\
\nt{VBD} &\rightarrow& \text{obtained} \mid \text{drew} \\
\nt{IN} &\rightarrow& \text{from} \mid \text{of} \\
\nt{CC} &\rightarrow& \text{or} \mid \text{and}.
\end{array}
\]

Deze grammatica bevat precies de soorten regels die nodig zijn voor de drie gegeven CoreNLP-bomen. Zinnen bestaan uit een subject-$\nt{NP}$ en een $\nt{VP}$. Een $\nt{VP}$ kan met of zonder $\nt{PP}$ voorkomen. Een $\nt{NP}$ kan een $\nt{PP}$ krijgen, en twee $\nt{NP}$'s kunnen met \lex{or} of \lex{and} gecoördineerd worden.

%------------------------------------------------------------
\item \textbf{Ambiguïteit volgens deze CFG}

Een zin is syntactisch ambigu als dezelfde zin volgens dezelfde grammatica meer dan één parse tree kan krijgen. Volgens de CFG uit vraag 1 zijn alle drie de zinnen ambigu.

\begin{enumerate}
  \item \lex{the architect obtained a sketchbook or a pencil from the landlord}

  Deze zin heeft volgens de CFG \textbf{3 mogelijke parse trees}. De ambiguïteit zit in de $\nt{PP}$ \lex{from the landlord} en in de structuur van de coördinatie met \lex{or}.

  De drie mogelijkheden zijn:
  \begin{enumerate}
    \item De $\nt{PP}$ hoort bij de hele $\nt{VP}$:
    \[
    [\nt{VP}\ \lex{obtained}\ [\nt{NP}\ \lex{a sketchbook or a pencil}]\ [\nt{PP}\ \lex{from the landlord}]].
    \]

    \item De $\nt{PP}$ hoort bij de hele object-$\nt{NP}$:
    \[
    [\nt{VP}\ \lex{obtained}\ [\nt{NP}\ [\nt{NP}\ \lex{a sketchbook or a pencil}]\ [\nt{PP}\ \lex{from the landlord}]]].
    \]

    \item De $\nt{PP}$ hoort alleen bij de tweede $\nt{NP}$:
    \[
    [\nt{VP}\ \lex{obtained}\ [\nt{NP}\ [\nt{NP}\ \lex{a sketchbook}]\ \lex{or}\ [\nt{NP}\ \lex{a pencil from the landlord}]]].
    \]
  \end{enumerate}

  Dit is dus zowel \nt{PP}-attachment ambiguity als coördinatie-/scope-ambiguïteit.

  \item \lex{the architect drew a sketch of a house and a garden}

  Deze zin heeft volgens de CFG \textbf{3 mogelijke parse trees}. De ambiguïteit zit in de $\nt{PP}$ \lex{of a house and a garden} en in de coördinatie met \lex{and}.

  De drie mogelijkheden zijn:
  \begin{enumerate}
    \item De $\nt{PP}$ hoort bij de $\nt{VP}$:
    \[
    [\nt{VP}\ \lex{drew}\ [\nt{NP}\ \lex{a sketch}]\ [\nt{PP}\ \lex{of a house and a garden}]].
    \]

    \item De $\nt{PP}$ hoort bij de $\nt{NP}$ \lex{a sketch}:
    \[
    [\nt{VP}\ \lex{drew}\ [\nt{NP}\ [\nt{NP}\ \lex{a sketch}]\ [\nt{PP}\ \lex{of a house and a garden}]]].
    \]

    \item De coördinatie zit hoger:
    \[
    [\nt{VP}\ \lex{drew}\ [\nt{NP}\ [\nt{NP}\ \lex{a sketch of a house}]\ \lex{and}\ [\nt{NP}\ \lex{a garden}]]].
    \]
  \end{enumerate}

  Dit is dus \nt{PP}-attachment ambiguity en coördinatie-/scope-ambiguïteit.

  \item \lex{the architect obtained the blueprint of the museum and a donation}

  Deze zin heeft volgens de CFG \textbf{3 mogelijke parse trees}. De ambiguïteit zit in de $\nt{PP}$ \lex{of the museum} en in de coördinatie met \lex{and}.

  De drie mogelijkheden zijn:
  \begin{enumerate}
    \item CoreNLP-achtige lezing:
    \[
    [\nt{VP}\ \lex{obtained}\ [\nt{NP}\ [\nt{NP}\ \lex{the blueprint}]\ [\nt{PP}\ \lex{of the museum and a donation}]]].
    \]

    \item Meest plausibele lezing:
    \[
    [\nt{VP}\ \lex{obtained}\ [\nt{NP}\ [\nt{NP}\ \lex{the blueprint of the museum}]\ \lex{and}\ [\nt{NP}\ \lex{a donation}]]].
    \]

    \item De $\nt{PP}$ hoort bij de $\nt{VP}$:
    \[
    [\nt{VP}\ \lex{obtained}\ [\nt{NP}\ \lex{the blueprint}]\ [\nt{PP}\ \lex{of the museum and a donation}]].
    \]
  \end{enumerate}

  Dit is dus \nt{PP}-attachment ambiguity en coördinatie-/scope-ambiguïteit.
\end{enumerate}

%------------------------------------------------------------
\item \textbf{Correcte bomen volgens de meest plausibele betekenis}

\textbf{Zin 1.} De meest plausibele structuur is dat \lex{from the landlord} bij het verkrijgen hoort, en dat \lex{a sketchbook or a pencil} samen het object is.

\begin{center}
\resizebox{0.98\textwidth}{!}{%
\Tree [.\nt{S}
        [.\nt{NP} [.\nt{DT} the ] [.\nt{NN} architect ] ]
        [.\nt{VP} [.\nt{VBD} obtained ]
             [.\nt{NP} [.\nt{NP} [.\nt{DT} a ] [.\nt{NN} sketchbook ] ]
                  [.\nt{CC} or ]
                  [.\nt{NP} [.\nt{DT} a ] [.\nt{NN} pencil ] ] ]
             [.\nt{PP} [.\nt{IN} from ]
                  [.\nt{NP} [.\nt{DT} the ] [.\nt{NN} landlord ] ] ] ] ]
}
\end{center}

\textbf{Zin 2.} De meest plausibele structuur is dat \lex{of a house and a garden} bij \lex{a sketch} hoort.

\begin{center}
\resizebox{0.98\textwidth}{!}{%
\Tree [.\nt{S}
        [.\nt{NP} [.\nt{DT} the ] [.\nt{NN} architect ] ]
        [.\nt{VP} [.\nt{VBD} drew ]
             [.\nt{NP} [.\nt{NP} [.\nt{DT} a ] [.\nt{NN} sketch ] ]
                  [.\nt{PP} [.\nt{IN} of ]
                       [.\nt{NP} [.\nt{NP} [.\nt{DT} a ] [.\nt{NN} house ] ]
                            [.\nt{CC} and ]
                            [.\nt{NP} [.\nt{DT} a ] [.\nt{NN} garden ] ] ] ] ] ] ]
}
\end{center}

\textbf{Zin 3.} De meest plausibele structuur is dat de architect twee dingen verkreeg: \lex{the blueprint of the museum} en \lex{a donation}.

\begin{center}
\resizebox{0.98\textwidth}{!}{%
\Tree [.\nt{S}
        [.\nt{NP} [.\nt{DT} the ] [.\nt{NN} architect ] ]
        [.\nt{VP} [.\nt{VBD} obtained ]
             [.\nt{NP} [.\nt{NP} [.\nt{NP} [.\nt{DT} the ] [.\nt{NN} blueprint ] ]
                       [.\nt{PP} [.\nt{IN} of ]
                            [.\nt{NP} [.\nt{DT} the ] [.\nt{NN} museum ] ] ] ]
                  [.\nt{CC} and ]
                  [.\nt{NP} [.\nt{DT} a ] [.\nt{NN} donation ] ] ] ] ]
}
\end{center}

%------------------------------------------------------------
\item \textbf{PCFG op basis van de correcte treebank uit vraag 3}

De correcte treebank bestaat uit de drie bomen uit vraag 3. De kans van een regel wordt geschat met relatieve frequentie:
\[
P(A\rightarrow \beta)=\frac{\#(A\rightarrow \beta)}{\#(A)}.
\]
Hier is $\#(A)$ het totale aantal regels in de treebank met linkerkant $A$.

\begin{center}
\small
\begin{longtable}{llll}
\toprule
Regel & Telling regel & Telling linkerkant & Kans \\
\midrule
\endhead
$\nt{S}\rightarrow \nt{NP}\ \nt{VP}$ & 3 & 3 & $1$ \\
$\nt{VP}\rightarrow \nt{VBD}\ \nt{NP}$ & 2 & 3 & $\frac{2}{3}$ \\
$\nt{VP}\rightarrow \nt{VBD}\ \nt{NP}\ \nt{PP}$ & 1 & 3 & $\frac{1}{3}$ \\
$\nt{NP}\rightarrow \nt{DT}\ \nt{NN}$ & 12 & 17 & $\frac{12}{17}$ \\
$\nt{NP}\rightarrow \nt{NP}\ \nt{PP}$ & 2 & 17 & $\frac{2}{17}$ \\
$\nt{NP}\rightarrow \nt{NP}\ \nt{CC}\ \nt{NP}$ & 3 & 17 & $\frac{3}{17}$ \\
$\nt{PP}\rightarrow \nt{IN}\ \nt{NP}$ & 3 & 3 & $1$ \\
\midrule
$\nt{DT}\rightarrow \text{the}$ & 6 & 12 & $\frac{1}{2}$ \\
$\nt{DT}\rightarrow \text{a}$ & 6 & 12 & $\frac{1}{2}$ \\
$\nt{NN}\rightarrow \text{architect}$ & 3 & 12 & $\frac{1}{4}$ \\
$\nt{NN}\rightarrow \text{sketchbook}$ & 1 & 12 & $\frac{1}{12}$ \\
$\nt{NN}\rightarrow \text{pencil}$ & 1 & 12 & $\frac{1}{12}$ \\
$\nt{NN}\rightarrow \text{landlord}$ & 1 & 12 & $\frac{1}{12}$ \\
$\nt{NN}\rightarrow \text{sketch}$ & 1 & 12 & $\frac{1}{12}$ \\
$\nt{NN}\rightarrow \text{house}$ & 1 & 12 & $\frac{1}{12}$ \\
$\nt{NN}\rightarrow \text{garden}$ & 1 & 12 & $\frac{1}{12}$ \\
$\nt{NN}\rightarrow \text{blueprint}$ & 1 & 12 & $\frac{1}{12}$ \\
$\nt{NN}\rightarrow \text{museum}$ & 1 & 12 & $\frac{1}{12}$ \\
$\nt{NN}\rightarrow \text{donation}$ & 1 & 12 & $\frac{1}{12}$ \\
$\nt{VBD}\rightarrow \text{obtained}$ & 2 & 3 & $\frac{2}{3}$ \\
$\nt{VBD}\rightarrow \text{drew}$ & 1 & 3 & $\frac{1}{3}$ \\
$\nt{IN}\rightarrow \text{from}$ & 1 & 3 & $\frac{1}{3}$ \\
$\nt{IN}\rightarrow \text{of}$ & 2 & 3 & $\frac{2}{3}$ \\
$\nt{CC}\rightarrow \text{or}$ & 1 & 3 & $\frac{1}{3}$ \\
$\nt{CC}\rightarrow \text{and}$ & 2 & 3 & $\frac{2}{3}$ \\
\bottomrule
\end{longtable}
\end{center}

Controle van de belangrijkste tellingen:
\begin{itemize}
  \item $\nt{VP}$ komt 3 keer voor: 2 keer als $\nt{VP}\rightarrow \nt{VBD}\ \nt{NP}$ en 1 keer als $\nt{VP}\rightarrow \nt{VBD}\ \nt{NP}\ \nt{PP}$.
  \item $\nt{NP}$ komt 17 keer voor: 12 keer als $\nt{NP}\rightarrow \nt{DT}\ \nt{NN}$, 2 keer als $\nt{NP}\rightarrow \nt{NP}\ \nt{PP}$, en 3 keer als $\nt{NP}\rightarrow \nt{NP}\ \nt{CC}\ \nt{NP}$.
  \item $\nt{PP}$ komt 3 keer voor, steeds als $\nt{PP}\rightarrow \nt{IN}\ \nt{NP}$.
\end{itemize}

%------------------------------------------------------------
\item \textbf{Kans van de CoreNLP-boom en de correcte boom voor zin 1}

De kans van een parse tree is het product van de kansen van alle regels die in die boom gebruikt worden.

Voor de CoreNLP-boom van zin 1 wordt de $\nt{PP}$ \lex{from the landlord} binnen de rechter conjunct-$\nt{NP}$ geplaatst:
\[
\lex{a sketchbook or [a pencil from the landlord]}.
\]

\[
\begin{aligned}
P(T_{\text{CoreNLP}})
&=1\cdot \left(\frac{12}{17}\right)^4\cdot \left(\frac{1}{2}\right)^4
\cdot \frac{1}{4}\cdot \left(\frac{1}{12}\right)^3\\
&\quad \cdot \frac{2}{3}\cdot \frac{2}{3}\cdot \frac{3}{17}\cdot \frac{1}{3}
\cdot \frac{2}{17}\cdot 1\cdot \frac{1}{3}\\
&=\frac{1}{434476242}\\
&\approx 0.000000002302.
\end{aligned}
\]

Voor de correcte boom van zin 1 hoort de $\nt{PP}$ \lex{from the landlord} bij de $\nt{VP}$, en is \lex{a sketchbook or a pencil} samen het object.

\[
\begin{aligned}
P(T_{\text{correct}})
&=1\cdot \left(\frac{12}{17}\right)^4\cdot \left(\frac{1}{2}\right)^4
\cdot \frac{1}{4}\cdot \left(\frac{1}{12}\right)^3\\
&\quad \cdot \frac{1}{3}\cdot \frac{2}{3}\cdot \frac{3}{17}\cdot \frac{1}{3}
\cdot 1\cdot \frac{1}{3}\\
&=\frac{1}{102229704}\\
&\approx 0.000000009782.
\end{aligned}
\]

De kansen zijn dus verschillend. De correcte boom heeft een hogere kans dan de CoreNLP-boom:
\[
P(T_{\text{correct}}) > P(T_{\text{CoreNLP}}).
\]
De correcte boom is ongeveer $4.25$ keer zo waarschijnlijk volgens deze PCFG.

%------------------------------------------------------------
\item \textbf{Evaluatie van CoreNLP tegenover de gold standard}

Ik gebruik labeled constituent evaluation. Dat betekent dat een constituent alleen correct telt als zowel het label als de woordspan hetzelfde zijn. Ik tel hier alleen constituent-nodes met labels $\nt{S}$, $\nt{NP}$, $\nt{VP}$ en $\nt{PP}$. Ik tel de woorden zelf en de PoS-labels $\nt{DT}$, $\nt{NN}$, $\nt{VBD}$, $\nt{IN}$ en $\nt{CC}$ niet mee. De root-node $\nt{S}$ tel ik dus wel mee als constituent.

De formules zijn:
\[
\begin{aligned}
\text{Precision}
&=\frac{\#\text{correcte CoreNLP-constituents}}{\#\text{CoreNLP-constituents}},\\
\text{Recall}
&=\frac{\#\text{correcte CoreNLP-constituents}}{\#\text{gold-standard constituents}}.
\end{aligned}
\]

\[
F_1 = \frac{2\cdot \text{Precision}\cdot \text{Recall}}{\text{Precision}+\text{Recall}}.
\]

\begin{center}
\small
\begin{tabular}{lllllll}
\toprule
Zin & CoreNLP & Gold & Correct & Precision & Recall & $F_1$ \\
\midrule
1 & 9 & 8 & 7 & $\frac{7}{9}\approx 0.778$ & $\frac{7}{8}=0.875$ & $\frac{14}{17}\approx 0.824$ \\
2 & 8 & 9 & 8 & $\frac{8}{8}=1.000$ & $\frac{8}{9}\approx 0.889$ & $\frac{16}{17}\approx 0.941$ \\
3 & 9 & 9 & 7 & $\frac{7}{9}\approx 0.778$ & $\frac{7}{9}\approx 0.778$ & $\frac{7}{9}\approx 0.778$ \\
\bottomrule
\end{tabular}
\end{center}

Uitwerking per zin:
\begin{enumerate}
  \item Zin 1: CoreNLP heeft 9 constituents, de gold standard heeft 8 constituents, en 7 daarvan komen overeen. Dus:
  \[
  P=\frac{7}{9},\qquad R=\frac{7}{8},\qquad F_1=\frac{2\cdot 7}{9+8}=\frac{14}{17}.
  \]

  \item Zin 2: CoreNLP heeft 8 constituents, de gold standard heeft 9 constituents, en 8 daarvan komen overeen. Dus:
  \[
  P=\frac{8}{8}=1,
  \qquad R=\frac{8}{9},
  \qquad F_1=\frac{2\cdot 8}{8+9}=\frac{16}{17}.
  \]

  \item Zin 3: CoreNLP heeft 9 constituents, de gold standard heeft 9 constituents, en 7 daarvan komen overeen. Dus:
  \[
  P=\frac{7}{9},\qquad R=\frac{7}{9},\qquad F_1=\frac{2\cdot 7}{9+9}=\frac{7}{9}.
  \]
\end{enumerate}

\end{enumerate}

%%%%%%%%%%%%%%%%%%%%%%%%%%%%%%%%%%%%
%% END
%%%%%%%%%%%%%%%%%%%%%%%%%%%%%%%%%%%%

\end{document}
